% Academic report shape: one skeleton for four kinds, every heading in French and English.
% Same parts and order as academic-report.md; the comment above each part says which kinds need it.
% Kinds: PR project report, IR internship report, MT master's thesis, TR technical report.
% Delete the parts your kind leaves out. Standard article class; compiles with pdflatex (run it twice).
\documentclass[11pt,a4paper]{article}
\usepackage[T1]{fontenc}
\usepackage[utf8]{inputenc}
\usepackage[french,english]{babel}

% A heading in both languages: \bi{French}{English}.
\newcommand{\bi}[2]{#1 \textperiodcentered{} #2}

% Title page, all kinds: the finding, then the school's items. State the school's length limit here if it sets one.
\title{Title that states the finding}
\author{Author\\ \small School and degree, course or host organisation\\ \small Supervisor: name}
\date{Date}

\begin{document}
\maketitle

% All kinds but PR, where it is optional: question, method, result and limit, about 200 words.
% The same text in French and English when the school asks, with at most five keywords.
\section*{\bi{Résumé}{Abstract}}
The question, the method, the result and its limit.

\noindent\textbf{Mots-clés · Keywords:} at most five terms.

% MT and TR; optional for PR and IR: the answer and what to do, in one page.
\section*{\bi{Synthèse}{Executive summary}}
The answer and what to do, for a reader who reads nothing else.

% MT and TR; PR and IR past 20 pages.
\renewcommand{\contentsname}{\bi{Sommaire}{Contents}}
\tableofcontents

% All kinds: the question, its scope and the plan. MT: end with the contribution in one sentence.
% PR and TR: the background the question needs, briefly, here rather than in a chapter.
\section{Introduction}
The question, its scope and the plan of the report.

% IR only: one page; the rest goes to an appendix.
\section{\bi{Contexte}{Host organisation}}
What the organisation does, and where the internship sat in it.

% PR and IR: the objective and the test that says it is met; IR: the learning objectives too.
\section{\bi{Mission et objectifs}{Mission and objectives}}
The objective, and how to tell it is met.

% MT only, as a chapter: open with a one-paragraph answer; end each section on what it implies for the question.
\section{\bi{État de l'art}{Literature and background}}
The answer first: what is known, what is missing, and where this thesis stands.

% Examiners often form their first impression here (Mullins 2002; Holbrook 2007, both of doctoral theses).
\subsection*{\bi{Champ et recherche}{Scope and search}}
What the chapter covers, and how the sources were found.

\subsection*{\bi{Thèmes}{Themes}}
One section per theme, not per source, each ending on what it implies for the question.

\subsection*{\bi{Le manque}{The gap}}
What no source settles, stated as a claim this thesis can test.

\subsection*{Position}
Where this thesis stands, and the method the gap calls for.

% All kinds: the approach, why this one, what it cannot show; one step per subsection, named by what it achieves.
\section{\bi{Méthode}{Method}}
\subsection{First step, named by what it achieves}
The figure of the process first, then how and why.

% All kinds: one subsection per result, its headline sentence first, the figure beside its text.
\section{\bi{Résultats}{Results}}
\subsection{First result, stated as a sentence}
The headline sentence, the number, then the figure.

% All kinds: what the results mean, how they compare, their limits; no claim beyond the data.
\section{Discussion}
What the results mean, and their limits.

% IR only: what was learnt, against the learning objectives of the mission section.
\section{\bi{Retour réflexif}{Reflection}}
What was learnt, objective by objective.

% All kinds: the answer again, standing alone. Nothing new.
\section{Conclusion}
The answer, so that a reader who opens the report here understands it.

% IR, MT and TR, only when the report proposes action.
\section{\bi{Recommandations et perspectives}{Recommendations and perspectives}}
What to do next, and why.

% PR and TR written by more than one author: who did and wrote what.
\section*{Contributions}
Who did and wrote what.

% All kinds, when the school or the course asks: whether and how AI tools were used, in its words.
\section*{\bi{Déclaration d'usage de l'IA}{Use of AI tools}}
Whether and how AI tools were used.

% Optional, IR and MT, at the back so that the answer comes first.
\section*{\bi{Remerciements}{Acknowledgements}}
Thanks.

% MT and TR: every term defined at first use, and listed here.
\section*{\bi{Glossaire}{Glossary}}
\begin{description}
  \item[Term] Its definition, as at its first use.
\end{description}

% All kinds, in the school's style; every entry with a DOI or a stable URL.
\renewcommand{\refname}{\bi{Bibliographie}{References}}
\begin{thebibliography}{9}
  \bibitem{placeholder} Author (Year). Title. Venue. DOI.
\end{thebibliography}

% Detail a reader can skip: the host organisation's longer presentation, raw data, code listings.
\appendix
\section{\bi{Annexes}{Appendices}}
Supporting material.

\end{document}
